%% =====================================================================
%% Thumb tabs (switched on and styled in settings/thumbs.tex)
%% A tab with the chapter number on the outer page edge, running into the bleed, so the
%% chapters can be found by flipping through the closed book.
%%  - Tabs step down the page, one position per chapter; appendices continue below the last
%%    chapter instead of starting again at the top.
%%  - The spacing adapts to the number of chapters and appendices (counted on the previous
%%    compile), so the tabs always fit between the top and bottom of the text block, on any
%%    page size. With few chapters, tabs are spaced 2 tab heights apart; with many, they move
%%    closer and become shorter.
%%  - Square tabs in a light grey with a black number ("thumb" and "thumbtext" in
%%    settings/thumbs.tex). Grey prints with black ink only; coloured tabs would make every page of
%%    the book a colour page at the printer.
%% \thumbsoff and \thumbson switch them off and on, e.g. around the references.
%% =====================================================================
\newlength\thumbheight      % tab height (maximum)                 (these three are set in
\newlength\thumbedge        % padding around the number             settings/thumbs.tex)
\newlength\thumbhspace      % how far the tab reaches into the page
\newif\ifthesisthumbs \newif\ifthesisthumbsdigital \thesisthumbsdigitaltrue
%% =====================================================================
%% Thumb tabs: a tab with the chapter number on the outer page edge, so the chapters can be found by
%% flipping through the closed book (the drawing is in template/thumb.tex).
%% =====================================================================
\thesisthumbstrue           % \thesisthumbsfalse for none at all
\thesisthumbsdigitaltrue    % also in the digital edition (on screen they mark the chapters at the page edge,
                            % as in the printed book); \thesisthumbsdigitalfalse for the print editions only

%% The tab and its number: light grey with a black number by default; for the darker look, anthracite
%% (or tudelft-black) with a white number. Any colour other than grey makes every page with a tab a
%% colour page at the printer.
\colorlet{thumb}{lightgrey}
\colorlet{thumbtext}{tudelft-black}
% \colorlet{thumb}{anthracite}\colorlet{thumbtext}{tudelft-white}

%% Size: the tab's height (at most; with many chapters the tabs get shorter so they all fit), the space
%% around its number, and how far it reaches into the page.
\setlength\thumbheight{24pt}
\setlength\thumbedge{4pt}
\setlength\thumbhspace{36pt}

\ifprintversion\else\ifthesisthumbsdigital\else\thesisthumbsfalse\fi\fi % (the digital edition)
\newlength\thumbvspace                                      % step between tab positions (computed)
\newlength\thesis@thumbtab                                  % actual tab height (computed)

\newif\ifthesis@thumbson \thesis@thumbsontrue
\newcommand{\thumbsoff}{\thesis@thumbsonfalse}
\newcommand{\thumbson}{\thesis@thumbsontrue}

%% Tab position = chapter number, plus the number of main chapters once in the appendices.
\newcounter{thesis@mainchapters}
\AddToHook{cmd/appendix/before}{\setcounter{thesis@mainchapters}{\value{chapter}}}
\newcommand{\thesis@thumbpos}{\the\numexpr\value{chapter}+\value{thesis@mainchapters}\relax}

%% Total number of tab positions, saved for the next compile.
\AtEndDocument{\immediate\write\@auxout{\gdef\string\thesis@thumbtotal{\thesis@thumbpos}}}
\AtBeginDocument{%
    \ifdefined\thesis@thumbtotal\else\gdef\thesis@thumbtotal{1}\fi
    \ifnum\thesis@thumbtotal<1 \gdef\thesis@thumbtotal{1}\fi
    %% step = text height / number of tabs, but at most 2 tab heights
    \setlength\thumbvspace{\dimexpr\textheight/\thesis@thumbtotal\relax}%
    \ifdim\thumbvspace>2\thumbheight \setlength\thumbvspace{2\thumbheight}\fi
    %% tab height = the maximum, or 80% of the step if the tabs are packed closely
    \setlength\thesis@thumbtab{\thumbheight}%
    \ifdim\thesis@thumbtab>0.8\thumbvspace \setlength\thesis@thumbtab{0.8\thumbvspace}\fi
}

%% The running head only marks the page (\lthumb on left pages, \rthumb on right pages); the tab itself
%% is drawn when the page is shipped out, in the page background, whose origin is the top left corner of
%% the sheet. Its position is then known exactly, without TikZ's "remember picture" (which stores
%% positions in the .aux file and makes latexmk compile once more whenever one shifts).
%% The book page sits inside the sheet (bleed, or A4 in the proof edition): \trimtop below the sheet's top,
%% \trimedge from its outer edge.
\usetikzlibrary{calc}
\gdef\thesis@thumbside{}
%% In a tagged PDF the tab is an artifact (decoration, not content), like the running head.
\ExplSyntaxOn
\cs_new_protected:Npn \__thesis_putthumb:
  { \put(0,0){\exp_after:wN \thesis@drawthumb \thesis@thumbside} }
\AddToHook{shipout/background}{
    \tl_if_empty:NF \thesis@thumbside
      {
        \tag_mc_end_push:
        \tag_mc_begin:n {artifact}
        \tag_suspend:n {thumb}
        \__thesis_putthumb:
        \tag_resume:n {thumb}
        \tag_mc_end:
        \tag_mc_begin_pop:n {}
        \tl_gclear:N \thesis@thumbside
      }
}
\ExplSyntaxOff

%% Draws the tab: from beyond the page edge (into the bleed) to \thumbhspace into the page, with the
%% number at its inner end. #1: the x position of the page edge on the sheet; #2: the direction into the
%% page, 1 or -1; #3: the side of the number its inner end is on.
\newcommand*{\thesis@drawthumb}[3]{%
    \begin{tikzpicture}[overlay]
        \coordinate (tab top) at (#1,-\trimtop-1in-\topmargin-\headheight-\headsep-\thesis@thumbpos\thumbvspace+\thumbvspace);
        \fill[fill=thumb] ($(tab top)-(#2*3*\thumbedge,0pt)$) rectangle ($(tab top)+(#2*\thumbhspace,-\thesis@thumbtab)$);
        \node at ($(tab top)+(#2*\thumbhspace,-0.5\thesis@thumbtab)$)[anchor=#3,inner sep=2\thumbedge]{%
            \headingfont\bfseries\color{thumbtext}\fontsize{0.75\thesis@thumbtab}{0.75\thesis@thumbtab}\selectfont\thechapter};
    \end{tikzpicture}}
\newcommand*{\thesis@thumb}[1]{%
    \ifthesisthumbs\ifthesis@thumbson\ifnum\value{chapter}>0
        \xdef\thesis@thumbside{#1}%
    \fi\fi\fi}
\newcommand*{\lthumb}{\thesis@thumb{{\the\trimedge}{1}{east}}}
\newcommand*{\rthumb}{\thesis@thumb{{\the\dimexpr\stockwidth-\trimedge\relax}{-1}{west}}}
